\begin{figure}
    \begin{lstlisting}[]
void spmv(row_ptrs: mref<? x idx>, // adj matrix CSR pointers
          col_idxs: mref<? x idx>, // adj matrix CSR indices
          nnz_vals: mref<? x f32>, // matrix nonzero values
          vec_vals: mref<? x f32>, // dense vector values
          out_vals: mref<? x f32>){ // output tensor
  // Go through all rows
  for (idx row_idx=0; row_idx<n_rows; row_idx++){
    idx ptr_beg = row_ptrs[row_idx];
    idx ptr_end = row_ptrs[row_idx+1];
    // Go through all nonzeros in a row
    for (idx ptr=ptr_beg; ptr<ptr_end; ptr++){
      f32 nnz_val = nnz_vals[ptr];
      idx col_idx = col_idxs[ptr];
      f32 vec_val = vec_vals[col_idx];
      out_vals[row_idx] += nnz_val * vec_val; }}}}
    \end{lstlisting}
    \vspace{-1em}
    \caption[Pseudocode of sparse matrix-vector multiplication.]{Pseudocode for the SpMV operation in \Cref{fig:spmv-lookup}. The outer loop iterates over CSR rows using \texttt{row\_ptrs}; the inner loop scans each row's nonzeros, loads the corresponding nonzero value and column index, and performs an indirect lookup into \texttt{vec\_vals}. The final statement accumulates the scaled dense-vector value into the output row. In the end, SpMV exposes data-dependent loop bounds and indirect load operations, which general-purpose CPUs can vectorize using gather operations.}
    \label{fig:spmv-code}
\end{figure}
